\hypertarget{sec:variables}{}%
\begin{center}
    {\Large\bfseries 6. VARIABLES}
\end{center}
\vspace{1em}

Las hipótesis del numeral 5 relacionan las siguientes variables, que se identifican, se clasifican según su función y se definen conceptualmente a partir del marco teórico. Su definición operacional (indicadores, escalas de medición, fuentes de datos y procedimiento de cálculo) se desarrollará en el diseño metodológico. Como el alcance es descriptivo con componente comparativo (numeral 4) y no se manipulan las caídas del servicio, la clasificación por función es analítica, no causal.

\begingroup
\small
\begin{longtable}{|p{3.5cm}|p{2.5cm}|p{6.5cm}|p{1.5cm}|}
\caption{Matriz de operacionalización y clasificación de variables} \label{tab:matriz_variables} \\
\hline
\textbf{Variable} & \textbf{Tipo} & \textbf{Definición Operacional / Conceptual} & \textbf{Hipótesis} \\ \hline
\endfirsthead

\multicolumn{4}{c}{{\thetable{} }} \\
\hline
\textbf{Variable} & \textbf{Tipo} & \textbf{Definición Operacional / Conceptual} & \textbf{Hipótesis} \\ \hline
\endhead



\hline
\endlastfoot

Nivel de dependencia funcional hacia la aplicación financiera digital & Dependiente (principal) & Grado en que el estudiante necesita una billetera digital o app bancaria móvil (Nequi, DaviPlata o Bancolombia) para realizar sus pagos cotidianos de alimentación, transporte y pagos académicos. Se manifiesta en la frecuencia de uso, la proporción de pagos que se hacen por este medio y la dificultad de reemplazarlo. & H1a, H1b \\ \hline
Aplicación financiera digital utilizada & Independiente (categórica) & Plataforma de pago móvil que el estudiante usa activamente: Nequi, DaviPlata o Bancolombia. Es el factor de comparación entre grupos. Un estudiante puede usar más de una, por lo que los grupos no son totalmente independientes. & H1b \\ \hline
Impacto operacional de la interrupción del servicio & Dependiente & Grado en que una caída, falla o latencia prolongada de la aplicación altera las actividades cotidianas del estudiante: aplazarlas, cancelarlas o modificarlas, y qué tan grave es la afectación percibida. & H1c \\ \hline
Mecanismos de contingencia ante la indisponibilidad & Dependiente & Estrategias con que el estudiante sigue con sus actividades cuando su billetera preferida no funciona: efectivo de reserva, cambio a otra plataforma, ayuda de pares o cancelación de la compra o el pago. & H1d \\ \hline
Tipo de actividad cotidiana afectada & Interviniente & Categoría de actividad que depende del pago digital: alimentación en cafeterías, transporte urbano, fotocopias e impresiones, arriendo u otras. La afectación varía según la urgencia y la posibilidad de sustituir el medio de pago. & H1a, H1c \\ \hline
Experiencia previa con caídas del servicio & Interviniente & Haber vivido o no una interrupción de la aplicación, y con qué frecuencia y duración. Condiciona la percepción del impacto y la preparación de planes de contingencia. Define además la subpoblación de H1c. & H1c, H1d \\ \hline
Hábitos de pago y disponibilidad de efectivo & Interviniente & Preferencia por pagos digitales frente al efectivo, y monto o hábito de mantener dinero físico. Modula tanto el nivel de dependencia como la capacidad de contingencia. & H1a, H1d \\ \hline
Características sociodemográficas y académicas & Interviniente / de control & Programa, semestre, edad, género y situación de vivienda (residente o no en zonas aledañas al campus), que pueden explicar diferencias en el uso y la dependencia. & H1a, H1b \\ \hline
\end{longtable}
\endgroup